\documentclass[border=5mm]{standalone}
\usepackage[siunitx, american]{circuitikz}

\begin{document}
\begin{circuitikz}[scale=1, transform shape, line width=0.6pt]
    
    % ---------------------------------------------------
    % ETAPA 1: Entrada y Par Darlington (Q1, Q2)
    % ---------------------------------------------------
    % Fuente de entrada y NodoInput
    \draw (0,0) node[ground]{} to[sV, l=$V_1$] (0,4) to[C, l=$10\mu$] (3,4) node[above]{B1};
          
    % Resistencia de base Q1
    \draw (3,4) to[R, l=$R_{BD}$] (3,0) node[ground]{};
    
    % Transistor Q1 (NPN)
    \draw (3,4) to[R, l=$h_{ie1}$, i=$i_{b1}$] (6,4) node[above right]{E1};
    \draw (6,7) node[ground, rotate=180]{} node[right]{C1} 
          to[cI, l=$\beta i_{b1}$] (6,4); 
    
    % Transistor Q2 (NPN)
    \draw (6,4) to[short] (8,4) node[above]{B2};
    \draw (8,4) to[R, l=$h_{ie2}$, i=$i_{b2}$] (11,4) node[above right]{E2};
    \draw (11,7) node[ground, rotate=180]{} node[right]{C2} to[cI, l=$\beta i_{b2}$] (11,4); 
    
    % Resistencia de emisor Q2
    \draw (11,4) to[R, l=$R_{ED}$] (11,0) node[ground]{};

    % ---------------------------------------------------
    % ETAPA 2: Emisor Común (Q3)
    % ---------------------------------------------------
    % Resistencias de polarización separadas (R1E y R2E)
    \draw (11,4) to[C, l=$10\mu$] (12.5,4);
    \draw (12.5,4) to[R, l=$R_{1E}$] (12.5,0) node[ground]{};
    \draw (12.5,4) to[short] (14,4);
    \draw (14,4) to[R, l=$R_{2E}$] (14,0) node[ground]{};
    
    % Transistor Q3 (NPN)
    \draw (14,4) to[short] (15.5,4) node[above]{B3};
    \draw (15.5,4) to[R, l=$h_{ie3}$, i=$i_{b3}$] (18,4) node[above right]{E3};
    \draw (18,4) to[R, l=$R_{EE}$] (18,0) node[ground]{}
    \draw (18,4) to[short] (20,4) to[C, l=$33\mu$] (20,0) node[ground]{};
    \draw (0,-6) node[above]{C3} to[cI, l=$\beta i_{b3}$] (0,-10) node[below right]{E3} node[ground]{};
      
% Resistencia de colector de Q3
\draw (0,-6) to[short] (1.5,-6) to[R, l=$R_{CE}$] (1.5,-10) node[ground]{};

% ---------------------------------------------------
% ETAPA 3: Par Sziklai y Carga (Q4, Q5, Salida)
% ---------------------------------------------------
% Resistencias de polarización separadas (R1S y R2S)
\draw (1.5,-6) to[short] (3,-6);
\draw (3,-6) to[R, l=$R_{1S}$] (3,-10) node[ground]{};
\draw (3,-6) to[C, l=$10\mu$] (4.5,-6);
\draw (4.5,-6) to[R, l=$R_{2S}$] (4.5,-10) node[ground]{};

% Transistor Q4 (NPN)
\draw (4.5,-6) to[short] (6,-6) node[above]{B4};
\draw (6,-6) to[R, l=$h_{ie4}$, i=$i_{b4}$] (9,-6) node[below right]{E4};
\draw (9,-3) node[left]{C4} to[cI, l=$\beta i_{b4}$] (9,-6);

% Transistor Q5 (PNP)
\draw (9,-3) node[above left]{B5} to[R, l=$h_{ie5}$, i<=$i_{b5}$] (12,-3) node[above right]{E5} node[ground, rotate=180]{};
\draw (12,-3) to[cI, l=$\beta i_{b5}$] (12,-6) node[below]{C5};
\draw (12,-6) to[short] (9,-6); 

% NodoOut y resistencia de emisor Sziklai
\draw (9,-6) to[short, -*] (11,-6);
\draw (11,-6) to[R, l=$R_{ES}$] (11,-10) node[ground]{};

% Red de Carga final (RO y RL en serie)
\draw (11,-6) to[short] (13,-6) to[R, l=$R_O$] (15,-6) node[above right]{Out} to[C, l=$10\mu$] (15,-8) to[R, l=$R_L$] (15,-10) node[ground]{};

\end{circuitikz}
\end{document}